\chapter*{Code Availability}
\addcontentsline{toc}{chapter}{Code Availability}

All code written for this thesis is available at
\todo{Insert the repository URL and state the licence. If the repository is private, say under
what conditions it can be made available, and give the commit hash or release tag corresponding
to the results reported here so that the version is unambiguous.}

\section*{Scripts}

The repository follows the two-block division used in this thesis: \texttt{block1\_cohort}
holds the work carried out on UK Biobank, \texttt{block2\_detection} the work on Ar-PlaqSegm1.
Each block separates scripts from results, and data locations are resolved centrally by
\texttt{common/paths.py} so that the scripts run independently of the working directory. The
pipeline is organised as standalone scripts rather than notebooks, so that each stage can be
re-run without manual intervention.

\begin{description}[leftmargin=2.8cm,style=nextline,nosep,itemsep=3pt]
  \item[\texttt{preprocess\_ukb\_files.py}] Extracts the UK Biobank image archives and writes the
        manifest consumed by the inference script.
  \item[\texttt{run\_plaque\_inference.py}] Applies the base detection model across the manifest
        and writes image-level detections and per-participant plaque counts.
  \item[\texttt{merge\_main\_axis.py}] Aggregates the long-axis selection into per-participant
        image counts.
  \item[\texttt{prepare\_yolo\_seg\_dataset.py}] Converts the Ar-PlaqSegm1 masks into YOLO
        polygon labels and writes the stratified training and validation split.
  \item[\texttt{prepare\_yolo\_det\_dataset.py}] Derives the corresponding bounding-box dataset
        used for the detection control.
  \item[\texttt{finetune\_yolo\_seg.py}] Fine-tunes the segmentation models; the four reported
        configurations differ only in the arguments given on the command line.
  \item[\texttt{finetune\_yolo\_det.py}] Trains the detection and no-copy-paste controls under an
        otherwise identical schedule.
  \item[\texttt{evaluate\_seg\_testset.py}] Evaluates a checkpoint on the held-out test partition
        and writes the confidence sweep and per-image results.
  \item[\texttt{compare\_variants.py}] Evaluates single models and pooled combinations on a
        common image set.
  \item[\texttt{threshold\_on\_val.py}] Implements the threshold protocol: selects the operating
        point on the validation split and reports the resulting test figures alongside the
        test-optimised ones.
  \item[\texttt{rescore\_instances.py}] Fits and evaluates the post-hoc instance re-scorer.
  \item[\texttt{make\_appendix\_gallery.py}] Renders the prediction sheets of
        Section~\ref{sec:appendix-gallery}, re-deriving the operating point from the validation
        split so that the panels and the reported figures cannot diverge.
  \item[\texttt{make\_pr\_curves.py}] Computes the instance-level precision-recall curves of
        Section~\ref{sec:appendix-training} from the validation split.
  \item[\texttt{make\_qualitative\_figure.py},\\\texttt{make\_echogenicity\_figure.py}]
        Render the image figures of Chapters~\ref{ch:introduction} and \ref{ch:results} from the
        dataset and the model checkpoints, so that every panel and every number printed on it is
        recomputed rather than transcribed.
  \item[\texttt{run\_cv.py}] Cross-validation over the full development set; provided for the
        continuation described in Section~\ref{sec:conclusion-outlook} and not used for the
        results reported here.
\end{description}

\section*{Software Environment}

Results were produced with Python 3.12 and the following package versions.

\begin{table}[H]
\centering
\caption[Software versions]{Package versions used for all reported results.}
\label{tab:software}
\begin{tabular}{ll}
\hline
Package & Version \\
\hline
\texttt{ultralytics} \cite{jocher2023ultralytics} & 8.3.248 \\
\texttt{torch} & 2.9.1 \\
\texttt{torchvision} & 0.24.1 \\
\texttt{opencv-python} & 4.12.0.88 \\
\texttt{numpy} & 2.2.6 \\
\texttt{pandas} & 2.3.3 \\
\texttt{scikit-learn} & 1.9.0 \\
\texttt{scipy} & 1.17.1 \\
\texttt{matplotlib} & 3.9.2 \\
\hline
\end{tabular}
\end{table}

Model training used the Metal Performance Shaders backend on Apple silicon; the cohort-scale
inference of Section~\ref{sec:reproduction} ran on cloud instances provided by the UK Biobank
Research Analysis Platform.
The gradient-boosted tree ensemble of Phase 1A used the XGBoost package as provided by the
Research Analysis Platform; the version string is not recoverable now that access has ended, and
is therefore not stated.

\section*{Data}

Ar-PlaqSegm1 is publicly available \cite{rulloni2025arplaqsegm1}. UK Biobank data cannot be
redistributed and are accessible only through the Research Analysis Platform under an approved
application.
